\documentclass[11pt,a4paper]{article}
\usepackage[utf8]{inputenc}
\usepackage[T1]{fontenc}
\usepackage[catalan]{babel}
\usepackage{lmodern}
\usepackage[sfdefault,lining]{FiraSans}
\usepackage[a4paper,top=2.3cm,bottom=2.2cm,left=1.8cm,right=1.8cm,headheight=1.3cm,headsep=0.4cm,footskip=1cm]{geometry}
\usepackage{amsmath,amssymb}
\usepackage{icomma}
\usepackage{graphicx}
\usepackage[table]{xcolor}
\usepackage{tikz}
\usepackage{fontawesome5}
\usepackage{ccicons}
\usepackage{fancyhdr}
\usepackage{lastpage}
\usepackage{tasks}
\usepackage{needspace}
\usepackage[most]{tcolorbox}
\usepackage{array}
\usepackage[hidelinks,pdftitle={Fitxa d'exercicis de la unitat 3: Nombres enters. 1r d'ESO},pdfsubject={Matemàtiques 1r ESO},pdfkeywords={nombres enters, recta numèrica, valor absolut, oposat, suma i resta, regla dels signes, operacions combinades, problemes}]{hyperref}

\definecolor{navy}{HTML}{1B2748}
\definecolor{accent}{HTML}{E8603A}
\definecolor{teal}{HTML}{1F9E90}
\definecolor{soft}{HTML}{EEF1F8}
\definecolor{muted}{HTML}{5D6685}

\setlength{\parindent}{0pt}
\setlength{\parskip}{3pt}
\settasks{label=\alph*), label-format=\bfseries\color{navy}, label-width=1.3em, item-indent=1.8em, column-sep=0.6em, after-item-skip=2pt, before-skip=3pt, after-skip=4pt}

% ---------- capçalera i peu ----------
\newcommand{\llicencia}{\href{https://creativecommons.org/licenses/by-nc-sa/4.0/deed.ca}{\ccbyncsa}\ \ Llicència Creative Commons Reconeixement-NoComercial-CompartirIgual 4.0 Internacional (CC BY-NC-SA 4.0)}
\pagestyle{fancy}
\fancyhf{}
\fancyhead[L]{\raisebox{-0.35\height}{\includegraphics[height=0.95cm]{logo-gva.png}}}
\fancyhead[C]{\small\color{muted}Fitxa d'exercicis · Unitat 3: Nombres enters · 1r d'ESO}
\fancyhead[R]{\raisebox{-0.35\height}{\includegraphics[height=0.95cm]{logo-vermellar.png}}}
\fancyfoot[L]{\scriptsize\color{muted}\llicencia}
\fancyfoot[R]{\small\color{muted}\thepage}
\renewcommand{\headrulewidth}{0.4pt}
\renewcommand{\headrule}{\color{muted!40}\hrule height \headrulewidth}
\fancypagestyle{primera}{\fancyhf{}\renewcommand{\headrulewidth}{0pt}
  \fancyfoot[L]{\scriptsize\color{muted}\llicencia}
  \fancyfoot[R]{\small\color{muted}\thepage}}

% ---------- nivells i exercicis ----------
\newcommand{\estrelles}[1]{{\color{accent}\ifcase#1\or\faStar\faStar[regular]\faStar[regular]\or\faStar\faStar\faStar[regular]\or\faStar\faStar\faStar\fi}}
\newcounter{exer}
\newcommand{\ex}[1]{\par\needspace{4\baselineskip}\medskip\stepcounter{exer}\noindent{\color{navy}\bfseries\theexer.}\ \ {\small\estrelles{#1}}\ \ }

\newcommand{\seccio}[2]{\par\needspace{10\baselineskip}\bigskip
  \begin{tcolorbox}[colback=navy,colframe=navy,coltext=white,arc=2mm,boxrule=0pt,left=3mm,right=3mm,top=1.2mm,bottom=1.2mm]
  {\large\bfseries #1\quad #2}
  \end{tcolorbox}\vspace{-2pt}}

\newcolumntype{C}[1]{>{\centering\arraybackslash}p{#1}}

% quadre de recordatori (només a la fitxa d'exercicis)
\newcommand{\recorda}[1]{\par\smallskip\begin{tcolorbox}[colback=accent!7,colframe=accent!55,arc=2mm,boxrule=0.5pt,left=2.5mm,right=2.5mm,top=1mm,bottom=1mm]\small{\color{accent}\faLightbulb[regular]}\ \textbf{Recorda:}\ #1\end{tcolorbox}}
% recta numèrica entera: \rectabuida{min}{max} i \recta{min}{max}{punts}
\newcommand{\rectabase}[2]{\draw[latex-latex,muted] (#1-0.8,0)--(#2+0.8,0);
  \foreach \i in {#1,...,#2}{\draw[muted] (\i,0.12)--(\i,-0.12) node[below=-1pt,font=\scriptsize,text=muted]{$\i$};}}
\newcommand{\rectabuida}[2]{\begin{tikzpicture}[x=0.78cm,baseline=0]\rectabase{#1}{#2}\end{tikzpicture}}
\newcommand{\recta}[3]{\begin{tikzpicture}[x=0.78cm,baseline=0]\rectabase{#1}{#2}
  \foreach \p in {#3}{\fill[accent] (\p,0) circle (0.11) node[above=1pt,font=\small\bfseries,text=navy]{$\p$};}\end{tikzpicture}}
\newcommand{\res}[1]{#1}

\begin{document}
\thispagestyle{primera}

\begin{minipage}[c]{0.25\textwidth}\includegraphics[width=\linewidth]{logo-gva.png}\end{minipage}%
\begin{minipage}[c]{0.5\textwidth}\centering
  {\color{navy}\Large\bfseries Fitxa d'exercicis de la unitat 3:\\[2pt] Nombres enters}\\[5pt]
  {\color{accent}\large\bfseries 1r d'ESO}
\end{minipage}%
\begin{minipage}[c]{0.25\textwidth}\hfill\includegraphics[width=0.85\linewidth]{logo-vermellar.png}\end{minipage}

\vspace{6pt}
{\color{navy}\hrule height 1.2pt}
\vspace{10pt}
Nom i cognoms: \rule{8.4cm}{0.4pt}\hfill Grup: \rule{1.6cm}{0.4pt}\hfill Data: \rule{2.2cm}{0.4pt}
\vspace{8pt}

\begin{tcolorbox}[colback=soft,colframe=soft,arc=2mm,left=3mm,right=3mm,top=1.5mm,bottom=1.5mm]
\small Exercicis graduats per nivells de dificultat:\quad \estrelles{1}\ bàsic \qquad \estrelles{2}\ mitjà \qquad \estrelles{3}\ avançat.\\[2pt] Resol-los al quadern mostrant tots els passos. Les rectes numèriques es poden completar directament en la fitxa. Fes els càlculs sense calculadora.
\end{tcolorbox}

% =====================================================
\seccio{1}{Els nombres enters}

\ex{1} Representa les situacions següents amb un nombre enter:\begin{tasks}(1)
  \task Un avió sobrevola la ciutat a 432 peus d'altura. \dotfill
  \task Aquesta nit el termòmetre ha marcat 6 °C baix zero. \dotfill
  \task El cotxe el tenim aparcat al tercer soterrani. \dotfill
  \task Sòcrates va nàixer l'any 470 abans de Crist. \dotfill
  \task M'he gastat tots els diners (ni en tinc ni en dec). \dotfill
  \task València està situada a 15 m sobre el nivell de la mar. \dotfill
  \task En un examen tipus test, cada resposta incorrecta resta 1 punt. \dotfill
\end{tasks}\recorda{Primer busca el \textbf{zero} (la referència): 0\,°C, la planta baixa, el nivell de la mar, no tindre ni deure diners\ldots{} El que està \emph{per damunt} és \textbf{positiu} i el que està \emph{per davall}, \textbf{negatiu}.}

\ex{1} Descriu tres situacions de la vida real en què s'utilitzen els nombres negatius.

\ex{1} Escriu una situació que es puga representar amb cada enter:\begin{tasks}(5)
  \task $-12$
  \task $+8$
  \task $-1500$
  \task $+2026$
  \task $-4$
\end{tasks}

\ex{1} Classifica aquests enters en \emph{positius}, \emph{negatius} i \emph{ni positius ni negatius}. Quins són també nombres naturals? \[-7\qquad 12\qquad 0\qquad +3\qquad {-25}\qquad 1\qquad {-1}\qquad 100\]

\ex{2} Escriu l'enter anterior i el següent de cada nombre:\begin{tasks}(5)
  \task $5$
  \task $-1$
  \task $0$
  \task $-10$
  \task $-99$
\end{tasks}

\ex{2} Escriu tots els enters compresos entre:\begin{tasks}(3)
  \task $-4$ i $3$
  \task $-10$ i $-6$
  \task $-2$ i $2$
\end{tasks}

\ex{3} Quants enters hi ha compresos entre $-50$ i $50$? I entre $-50$ i $0$?


% =====================================================
\seccio{2}{Recta numèrica, ordre i comparació}

\ex{1} Representa en la recta numèrica els nombres: \ $-2,\ 0,\ 4,\ {-1},\ 8,\ {-3},\ {-7}$ \ i \ $1$.\par\smallskip{\centering\rectabuida{-9}{9}\par}\recorda{Els \textbf{positius} van a la \textbf{dreta} del zero i els \textbf{negatius}, a l'\textbf{esquerra}. Com més a la dreta, més gran és el nombre.}

\ex{1} Representa en la recta numèrica aquests nombres i ordena'ls de menor a major: \[-7\qquad 3\qquad 1\qquad {-4}\qquad 6\qquad {-5}\qquad {-2}\qquad 0\]{\centering\rectabuida{-9}{9}\par}

\ex{1} Ordena els nombres següents de major a menor: \ $-7,\ 0,\ 1,\ {-1},\ 8,\ {-4},\ {-9},\ 3$.

\ex{1} Compara les parelles de nombres amb el signe $<$, $>$ o $=$:\begin{tasks}(4)
  \task $-6\ \square\ 6$
  \task $0\ \square\ {-1}$
  \task $10\ \square\ {-1}$
  \task $8\ \square\ 3$
  \task $-8\ \square\ {-9}$
  \task $1\ \square\ {-2}$
  \task $-15\ \square\ {-51}$
  \task $-100\ \square\ 1$
\end{tasks}

\ex{2} Escriu un enter comprès entre:\begin{tasks}(4)
  \task $-5$ i $-2$
  \task $-1$ i $1$
  \task $-13$ i $-11$
  \task $-100$ i $-98$
\end{tasks}

\ex{2} Quin enter està\ldots\begin{tasks}(2)
  \task 5 unitats a l'esquerra de $2$?
  \task 7 unitats a la dreta de $-4$?
  \task 3 unitats a l'esquerra de $-6$?
  \task a 4 unitats del $0$? (n'hi ha dos)
\end{tasks}

\ex{3} Vertader o fals? Justifica-ho:\begin{tasks}(1)
  \task Tot nombre negatiu és menor que zero.
  \task $-20>-2$.
  \task Entre $-3$ i $3$ hi ha 5 enters.
  \task El major dels nombres negatius és $-1$.
  \task Si $a<b$, aleshores $-a>-b$.
\end{tasks}


% =====================================================
\seccio{3}{Valor absolut i oposat}

\ex{1} Calcula el valor absolut dels nombres següents:\begin{tasks}(6)
  \task $\lvert -5\rvert $
  \task $\lvert -1\rvert $
  \task $\lvert 0\rvert $
  \task $\lvert 8\rvert $
  \task $\lvert -23\rvert $
  \task $\lvert 7\rvert $
\end{tasks}\recorda{El \textbf{valor absolut} d'un nombre és la \textbf{distància} d'aquest nombre al 0: \ $\lvert -5\rvert =\lvert 5\rvert =5$. Mai no és negatiu.}

\ex{1} Calcula l'oposat dels nombres següents:\begin{tasks}(6)
  \task $\text{op}(-8)$
  \task $\text{op}(-3)$
  \task $\text{op}(0)$
  \task $\text{op}(4)$
  \task $\text{op}(-100)$
  \task $\text{op}(15)$
\end{tasks}\recorda{L'\textbf{oposat} d'un nombre és el que està a la \textbf{mateixa distància} del zero, però a l'altre costat: canviem el signe. \ $\text{op}(-8)=8$, \ $\text{op}(4)=-4$.}

\ex{1} Dos nombres oposats es troben a una distància de $36$ unitats l'un de l'altre. Quins nombres són?

\ex{2} Ordena de menor a major els nombres següents, fent abans els càlculs: \[3\qquad \text{op}(4)\qquad \lvert -6\rvert \qquad 1\qquad \text{op}(5)\qquad {-2}\qquad \lvert 7\rvert \qquad \lvert -2\rvert \qquad 0\qquad 5\]

\ex{2} Compara les parelles de nombres:\begin{tasks}(3)
  \task $\lvert -3\rvert \ \square\ 3$
  \task $\text{op}(-2)\ \square\ {-1}$
  \task $\lvert -2\rvert \ \square\ \lvert -1\rvert $
  \task $\text{op}(\text{op}(-1))\ \square\ 0$
  \task $\lvert -4\rvert \ \square\ \text{op}(2)$
  \task $-8\ \square\ \text{op}(8)$
\end{tasks}

\ex{3} Completa els buits amb el nombre o els nombres que corresponguen:\begin{tasks}(3)
  \task $\text{op}(\text{op}(\square))=-7$
  \task $\lvert \square\rvert =10$
  \task $\text{op}(\lvert \square\rvert )=-2$
  \task $\lvert \text{op}(-8)\rvert =\square$
  \task $\text{op}(\lvert -4\rvert )=\square$
  \task $\lvert \square\rvert =\text{op}(5)$
\end{tasks}

\ex{3} Quins enters tenen valor absolut menor que $4$? I quins tenen valor absolut menor o igual que $2$?


% =====================================================
\seccio{4}{Suma i resta d'enters}

\ex{1} Calcula:\begin{tasks}(5)
  \task $5-7$
  \task $-8-9$
  \task $-3+7$
  \task $4-3$
  \task $-5+1$
  \task $10-14$
  \task $-7+5$
  \task $9-13$
  \task $3+10$
  \task $9-4$
  \task $-5-5$
  \task $-4+4$
  \task $-3-2$
  \task $8+12$
  \task $-5+15$
\end{tasks}\recorda{\textbf{Mateix signe}: se sumen i es posa el signe comú ($-3-5=-8$). \ \textbf{Signes diferents}: se resten (gran menys xicotet) i es posa el signe del que té major valor absolut ($-7+4=-3$).}

\ex{1} Calcula:\begin{tasks}(6)
  \task $10-3-5$
  \task $15-9-6$
  \task $5-8+4$
  \task $9-3+5$
  \task $-2+2+7$
  \task $-10+8+6$
  \task $-10-3+8$
  \task $-4-3-2$
  \task $-1-5-7$
  \task $9-2-3$
  \task $12-4-6$
  \task $3-7+4$
  \task $5-9+8$
  \task $-13+6+4$
  \task $-2+10-15$
  \task $-11-4+8$
  \task $-5-3-4$
  \task $-8+5+6$
\end{tasks}

\ex{2} Calcula:\begin{tasks}(4)
  \task $5+7-2-4$
  \task $2-6+4-9$
  \task $9-6-7+2$
  \task $-4-5+3+8$
  \task $-8+2-7+6$
  \task $-1+5+6-7$
  \task $4-1-2-11$
  \task $-2+6-8-1$
  \task $-1+2+17-12$
  \task $-24+15+13-8$
  \task $-18-10+12+1$
  \task $11-15-13-2$
\end{tasks}\recorda{Pots agrupar: suma tots els positius, suma tots els negatius i resta: \ $5-9+8-2=(5+8)-(9+2)=13-11=2$.}

\ex{2} Calcula de dues maneres, d'esquerra a dreta i agrupant positius i negatius:\begin{tasks}(2)
  \task $12-4-6+9-15$
  \task $-7+3-8+10-2+5$
\end{tasks}

\ex{3} Troba el nombre que falta:\begin{tasks}(3)
  \task $-3+\square=5$
  \task $\square-7=-2$
  \task $4+\square=-6$
  \task $\square+9=0$
  \task $-8-\square=-3$
  \task $\square-(-4)=1$
\end{tasks}


% =====================================================
\seccio{5}{Parèntesis i claudàtors}

\ex{1} Calcula, eliminant prèviament els parèntesis:\begin{tasks}(3)
  \task $(-6)+(-3)$
  \task $(+7)+(-4)$
  \task $(-8)+(+1)$
  \task $(-3)+(-7)$
  \task $(+12)+(+5)$
  \task $(+3)+(-5)$
  \task $(+3)+(+4)$
  \task $(+4)+(-6)$
  \task $(+7)+(-8)$
  \task $(-5)+(+4)$
  \task $(-9)+(+3)$
  \task $(+1)+(-1)$
  \task $(-7)+(+7)$
  \task $(-12)+(-3)$
  \task $(+7)+(+5)$
  \task $(-4)+(-10)$
  \task $(+6)+(-12)$
  \task $(-9)+(+15)$
\end{tasks}\recorda{Regla dels signes per a llevar parèntesis: \ $+(+\ \to\ +$ \qquad $+(-\ \to\ -$ \qquad $-(+\ \to\ -$ \qquad $-(-\ \to\ +$}

\ex{2} Calcula, eliminant prèviament els parèntesis:\begin{tasks}(2)
  \task $3-(-1)$
  \task $-8+(-17)$
  \task $-4-(-7)$
  \task $-(-5)+3$
  \task $8+(-14)$
  \task $(-2)-(+3)+(-6)$
  \task $-(-5)+(-4)-(-10)$
  \task $-15-(-3)+7$
  \task $(-8)+(-3)-(-17)+(+1)$
  \task $-2-(-3)+7-8$
  \task $-(-3)+(-1)-(+3)-(-8)$
  \task $-3+(-6)-4+1$
  \task $12-25+3-8-(-6)$
  \task $-7+4-(-8)+(-3)$
  \task $-(+4)+(-8)-(-6)+4-1$
  \task $-8+(-3)+4-7-5$
  \task $-(-6)+(-1)-7+(-3)$
  \task $15-(+3)-6+3-12$
\end{tasks}

\ex{2} Calcula, resolent prèviament els parèntesis:\begin{tasks}(4)
  \task $12+(+3-5)$
  \task $14+(+12-10)$
  \task $6-(9-7)$
  \task $15-(2-9)$
  \task $11-(-6+3)$
  \task $10-(-7-5)$
  \task $13+(-8+2)$
  \task $17+(-5-9)$
  \task $1+(-3+8)$
  \task $9-(2-1)$
  \task $-6-(2-7)$
  \task $11+(-1-3)$
  \task $6+(6+3)$
  \task $-5-(-2-8)$
  \task $3+(-8+7)$
  \task $15+(-6-8)$
\end{tasks}\recorda{Resol primer l'operació de dins del parèntesi, o bé lleva'l: si davant hi ha un $-$, canvia el signe de \textbf{tots} els termes de dins: \ $6-(9-7)=6-9+7=4$.}

\ex{2} Calcula, resolent prèviament els parèntesis:\begin{tasks}(3)
  \task $(7-4)+(9-5)$
  \task $(2+6)+(5-8)$
  \task $(5-9)+(2-12)$
  \task $(7+3)-(5+4)$
  \task $(8-12)-(2-5)$
  \task $(10-7)-(-2-6)$
  \task $-(8+4)+(5-9)$
  \task $-(6-2)-(7-9)$
  \task $(-8-2)+(-1-5)$
  \task $(7-1)-(-5+11)$
  \task $(15-8)+(-6-12)$
  \task $(-2+3)-(8-4)$
  \task $(17-12)-(-7-8)$
  \task $(-10+3)-(2-9)$
  \task $-(6-4)+(8+1)$
  \task $-(-7-1)-(2-10)$
\end{tasks}

\ex{3} Calcula, començant pels parèntesis de dins:\begin{tasks}(3)
  \task $6+[5+(7+2)]$
  \task $8+[4-(3+5)]$
  \task $10-[6+(2+7)]$
  \task $15-[2-(6-10)]$
  \task $15-[10-(8+4)]$
  \task $12-[7-(2-10)]$
  \task $(-6)+[5+(2-12)]$
  \task $(-7)-[3-(4-9)]$
  \task $(2-10)+[5-(8+2)]$
  \task $(12-3)-[1-(2-6)]$
  \task $[9-(+5)]+[7+(-10)]$
  \task $[10-(-2)]-[5-(+12)]$
  \task $[8-(6+4)]-(5-7)$
  \task $[1+(6-9)]-(8-12)$
  \task $2-[6-(-1-21)]$
  \task $-1+[-5-(2-13)]$
\end{tasks}


% =====================================================
\seccio{6}{Multiplicació i divisió}

\ex{1} Resol les multiplicacions següents:\begin{tasks}(3)
  \task $(-4)\cdot (+7)$
  \task $(+2)\cdot (+6)$
  \task $(-5)\cdot (-7)$
  \task $(+3)\cdot (-8)$
  \task $(-9)\cdot (-3)$
  \task $(-6)\cdot (+4)$
  \task $3\cdot (-2)$
  \task $4\cdot (+5)$
  \task $8\cdot (-6)$
  \task $-5\cdot (+3)$
  \task $-2\cdot (-4)$
  \task $-6\cdot (+3)$
  \task $(+3)\cdot (-5)\cdot (+2)$
  \task $(-4)\cdot (-1)\cdot (+6)$
  \task $(-2)\cdot (-7)\cdot (-2)$
  \task $(+5)\cdot (-4)\cdot (-3)$
  \task $(-8)\cdot 4\cdot (-6)$
  \task $5\cdot 3\cdot (-9)\cdot 4$
\end{tasks}\recorda{Regla dels signes: \ $+\cdot+\ \to\ +$ \qquad $+\cdot-\ \to\ -$ \qquad $-\cdot+\ \to\ -$ \qquad $-\cdot-\ \to\ +$. \ Per a la divisió és igual.}

\ex{1} Calcula:\begin{tasks}(3)
  \task $(-8):(+2)$
  \task $(+20):(-10)$
  \task $(-12):(-4)$
  \task $(-4):(+2)$
  \task $(+21):(-7)$
  \task $(-12):(+6)$
  \task $(-15):(-3)$
  \task $(+32):(+8)$
  \task $(-36):(+9)$
  \task $(+42):(-7)$
  \task $(-48):(-8)$
  \task $(+54):(+6)$
  \task $[(+80):(-8)]:(-5)$
  \task $[(-70):(-2)]:(-7)$
  \task $(+50):[(-30):(+6)]$
  \task $(-40):[(+24):(+3)]$
\end{tasks}

\ex{2} Troba el nombre que falta:\begin{tasks}(3)
  \task $(-6)\cdot(\ \square\ )=-18$
  \task $(\ \square\ )\cdot(-3)=-24$
  \task $(\ \square\ )\cdot(-5)=+35$
  \task $(+15)\cdot(\ \square\ )=+60$
  \task $(-8)\cdot(\ \square\ )=40$
  \task $(\ \square\ )\cdot7=-77$
\end{tasks}

\ex{2} Calcula (les multiplicacions i divisions seguides es fan d'esquerra a dreta):\begin{tasks}(4)
  \task $[(+6)\cdot (-4)]:(-3)$
  \task $[(-15)\cdot (-2)]:(+6)$
  \task $(-5)\cdot [(+12):(-3)]$
  \task $[(-5)\cdot (+12)]:(-3)$
  \task $24:(-2)\cdot (-3)$
  \task $(-40)\cdot [(-10):(-5)]$
  \task $(-70)\cdot (-5):(+2)$
  \task $(+100):(+4)\cdot (-2)$
\end{tasks}

\ex{3} Troba dos nombres enters que complisquen, en cada cas:\begin{tasks}(1)
  \task El producte és $-12$ i la suma és $1$.
  \task El producte és $18$ i la suma és $-9$.
  \task El producte és $-20$ i la suma és $-1$.
\end{tasks}

\ex{3} Sense fer els càlculs, digues quin signe tindrà el resultat. Després, comprova-ho:\begin{tasks}(1)
  \task $(-1)\cdot (-2)\cdot (-3)\cdot (-4)$
  \task $(-5)\cdot (+2)\cdot (-1)\cdot (-3)\cdot (+4)$
  \task El producte de set factors iguals a $-1$.
\end{tasks}


% =====================================================
\seccio{7}{Operacions combinades}

\ex{1} Calcula:\begin{tasks}(2)
  \task $(-8)\cdot (+2)+(-5)\cdot (-3)$
  \task $(+40):(-8)-(-30):(+6)$
  \task $(-2)\cdot (-9)+(-24):(-3)-(-6)\cdot 4$
  \task $(+27):(-3)-(+3)\cdot (-5)-(-6)\cdot (-2)$
  \task $(-3)\cdot [(-2)+(-4)]$
  \task $(+4)\cdot [(-5)+(+2)]$
  \task $(+6):[(+5)-(+7)]$
  \task $(-20):[(-6)-(-2)]$
\end{tasks}\recorda{Ordre de les operacions: \ 1) parèntesis i claudàtors (de dins cap a fora); \ 2) multiplicacions i divisions, d'esquerra a dreta; \ 3) sumes i restes.}

\ex{2} Calcula:\begin{tasks}(2)
  \task $[(-8)+(+7)]\cdot (-3)$
  \task $[(-9)+(-3)]:(+6)$
  \task $19-(-3)\cdot [5-(+8)]$
  \task $12+(-5)\cdot [8+(-9)]$
  \task $12-[13-(-7)]:(-5)$
  \task $10-(+20):[7+(-3)]$
  \task $(-2)\cdot (5-7)-(-3)\cdot (8-6)$
  \task $(9-6)\cdot (-2)+(13+3):(-4)$
\end{tasks}

\ex{3} Calcula:\begin{tasks}(2)
  \task $-2\cdot [3-(-4+6)\cdot (-2)]+15:(-5)$
  \task $[(-3)\cdot 4-(-2)\cdot 5]:(-2)$
  \task $18:[-7-(-4)]-(-2)\cdot [5-3\cdot (-1)]$
  \task $(-4)\cdot [6-2\cdot (3-7)]-[-10:(-5)]$
\end{tasks}

\ex{3} Posa parèntesis on calga perquè les igualtats siguen certes:\begin{tasks}(3)
  \task $6-4\cdot2-5=-6$
  \task $-3+5\cdot2-4=0$
  \task $10-2\cdot3+1=2$
\end{tasks}


% =====================================================
\seccio{8}{Problemes amb enters}

\ex{1} Quina diferència de temperatura suporta una persona que passa de la cambra de conservació de les verdures, que es troba a $4$\,°C, a la del peix congelat, que està a $-18$\,°C? I si passa de la cambra del peix a la de les verdures?\recorda{Escriu cada dada com un enter \textbf{amb el seu signe} (\emph{puja, guanya, sobre zero} $\to +$; \ \emph{baixa, perd, deu, baix zero} $\to -$). Per a calcular una \textbf{diferència} o variació: \ \textbf{valor final} $-$ \textbf{valor inicial}.}

\ex{1} La Mònica ix en l'ascensor des de la planta zero del seu edifici. L'ascensor puja 5 plantes, després baixa 3, puja 5, baixa 8, puja 10, puja 5 i baixa 6. En quina planta està?

\ex{1} En Joan li deu 14 euros a la seua germana. Si li paga 9 euros, quant li deu?

\ex{1} En una estació d'esquí, el termòmetre marcava $14$\,° baix zero a les 8 del matí; al migdia la temperatura havia pujat 10 graus i a les 19.00 havia baixat 5 graus respecte al migdia. Quina temperatura feia a les 19.00?

\ex{1} Si estava bussejant a una profunditat de 2 metres i ara estic sobre un penya-segat a 8 m d'altura, quina distància hi ha entre els dos llocs?

\ex{1} El primer mes de l'any la gasolina va baixar 5 cèntims, el segon va pujar 3 cèntims, el tercer va pujar 8 cèntims i el quart va baixar 2 cèntims. Expressa-ho amb nombres enters i calcula la variació total del preu en aquest temps.

\ex{2} Un dia d'hivern, la temperatura a les sis del matí era de $3$\,° baix zero. Al migdia havia pujat 9 graus, però a les dotze de la nit el termòmetre marcava $1$\,° baix zero. Quina diferència de temperatura hi ha hagut entre el migdia i la mitjanit?

\ex{2} Euclides va ser un matemàtic que va viure 60 anys i va morir l'any 265 a.\,C. En quin any va nàixer?

\ex{2} El Mulhacén, el pic més alt de la península Ibèrica, fa $3478$ metres d'altura. El Sistema del Trave, un dels avencs més profunds del món, arriba a $1441$ metres de profunditat. Troba la diferència d'altitud entre els dos punts.

\ex{2} En Pau tenia 850 euros a la seua cartilla d'estalvis. Ha ingressat 250 euros cada mes durant els 5 últims mesos i després en va traure 2300 per a pagar el fuster. Quin saldo li queda?

\ex{3} Arquimedes va nàixer l'any 287 a.\,C. i va morir l'any 212 a.\,C. Quants anys va viure?

\ex{3} Un submarí es troba a 120 m de profunditat. Puja 45 m i després baixa el triple del que ha pujat. A quina profunditat es troba ara?

\ex{3} En un examen tipus test, cada resposta correcta suma 3 punts, cada resposta incorrecta en resta 2 i cada pregunta en blanc val 0. La Laia n'ha encertat 11, n'ha fallat 6 i n'ha deixat 3 en blanc. Quina puntuació ha obtingut? Quantes n'hauria d'haver encertat, com a mínim, per a no tindre una puntuació negativa si n'haguera fallat 9 i n'haguera deixat 3 en blanc?

\vfill
\begin{center}\color{muted}\small\faCheckCircle[regular]\ \ Ànim! Si et costa algun exercici, repassa els apunts interactius de la unitat.\end{center}

\end{document}
